\documentclass[11pt,a4paper]{article}
\usepackage[utf8]{inputenc}
\usepackage[spanish]{babel}
\usepackage[T1]{fontenc}
\usepackage{geometry}
\geometry{margin=2cm}
\usepackage{xcolor}
\usepackage{listings}
\usepackage{tcolorbox}
\tcbuselibrary{most}
\usepackage{booktabs}
\usepackage{array}
\usepackage{hyperref}
\usepackage{fancyhdr}
\usepackage{titlesec}
\usepackage{multicol}
\usepackage{tikz}
\usetikzlibrary{arrows.meta,positioning,shapes.geometric,fit}

% Colores
\definecolor{tokunavy}{RGB}{20,55,100}
\definecolor{tokucyan}{RGB}{0,160,180}
\definecolor{codebg}{RGB}{245,245,245}
\definecolor{alertred}{RGB}{220,50,50}
\definecolor{tipgreen}{RGB}{39,174,96}
\definecolor{warnambel}{RGB}{230,126,34}
\definecolor{darkgray}{RGB}{50,50,50}
\definecolor{clpurple}{RGB}{130,60,160}
\definecolor{mxred}{RGB}{190,30,45}
\definecolor{brgreen}{RGB}{0,130,80}

\lstdefinestyle{python}{
  language=Python,
  backgroundcolor=\color{codebg},
  basicstyle=\ttfamily\small,
  keywordstyle=\color{tokunavy}\bfseries,
  stringstyle=\color{teal},
  commentstyle=\color{gray}\itshape,
  numbers=left, numberstyle=\tiny\color{gray},
  frame=single, framerule=0.4pt, rulecolor=\color{tokucyan!40},
  breaklines=true, tabsize=2, showstringspaces=false
}

% Soporte UTF-8 en listings (español)
\lstset{
  literate=
    {á}{{\'{a}}}1 {é}{{\'{e}}}1 {í}{{\'{i}}}1 {ó}{{\'{o}}}1 {ú}{{\'{u}}}1
    {Á}{{\'{A}}}1 {É}{{\'{E}}}1 {Í}{{\'{I}}}1 {Ó}{{\'{O}}}1 {Ú}{{\'{U}}}1
    {ñ}{{\~{n}}}1 {Ñ}{{\~{N}}}1 {ü}{{\"u}}1
    {¿}{{?`}}1 {¡}{{!`}}1
}

\newtcolorbox{tipbox}{colback=tipgreen!10,colframe=tipgreen,title=\textbf{Tip},fonttitle=\small\bfseries}
\newtcolorbox{alertbox}{colback=alertred!10,colframe=alertred,title=\textbf{Ojo},fonttitle=\small\bfseries}
\newtcolorbox{dominiobox}[1]{colback=tokunavy!5,colframe=tokunavy,title=\textbf{#1},fonttitle=\small\bfseries}
\newtcolorbox{clbox}{colback=clpurple!8,colframe=clpurple,title=\textbf{Chile},fonttitle=\small\bfseries}
\newtcolorbox{mxbox}{colback=mxred!8,colframe=mxred,title=\textbf{México},fonttitle=\small\bfseries}
\newtcolorbox{brbox}{colback=brgreen!8,colframe=brgreen,title=\textbf{Brasil},fonttitle=\small\bfseries}

\pagestyle{fancy}
\fancyhf{}
\lhead{\textcolor{tokunavy}{\textbf{Dominio: Pagos Recurrentes}}}
\rhead{\textcolor{gray}{Preparación Técnica — Toku}}
\cfoot{\thepage}

\titleformat{\section}{\large\bfseries\color{tokunavy}}{}{0em}{}[\titlerule]
\titleformat{\subsection}{\normalsize\bfseries\color{darkgray}}{}{0em}{}

\begin{document}

\begin{titlepage}
  \centering
  \vspace*{2cm}
  {\Huge\bfseries\color{tokunavy} Dominio: Pagos Recurrentes\par}
  \vspace{0.5cm}
  {\large\color{gray} El negocio de Toku explicado técnicamente\par}
  \vspace{0.3cm}
  \textcolor{tokucyan}{\rule{8cm}{2pt}}
  \vspace{1.5cm}

  \begin{tcolorbox}[colback=tokunavy!5,colframe=tokunavy,width=11cm]
    \centering
    \textbf{Por qué importa:} No puedes diseñar un sistema que no entiendes.\\[4pt]
    Toku procesa millones de cobros recurrentes en 3 países.\\[4pt]
    Cada país tiene su propio sistema bancario y regulación.
  \end{tcolorbox}

  \vspace{1.5cm}
  \begin{tabular}{ll}
    \textcolor{tokunavy}{\textbullet} & El modelo de negocio de Toku\\
    \textcolor{tokunavy}{\textbullet} & Débito directo vs tarjeta\\
    \textcolor{tokunavy}{\textbullet} & PAC/PAT (Chile), SPEI/CLABE (México), PIX (Brasil)\\
    \textcolor{tokunavy}{\textbullet} & Flujos de pago completos\\
    \textcolor{tokunavy}{\textbullet} & Idempotencia y reintentos\\
    \textcolor{tokunavy}{\textbullet} & La API de Toku desde adentro\\
  \end{tabular}
  \vfill
  {\small\color{gray} Marzo 2026 — Estudio Personal\par}
\end{titlepage}

\tableofcontents
\newpage

% ─── SECCIÓN 1: QUÉ HACE TOKU ────────────────────────────────────────────────
\section{¿Qué hace Toku exactamente?}

\begin{dominiobox}{Definición técnica}
Toku es una \textbf{plataforma de automatización de cobranza recurrente} que permite a
organizaciones cobrar directamente desde cuentas bancarias o tarjetas de sus clientes,
de forma automática y recurrente, sin que el pagador tenga que hacer nada.
\end{dominiobox}

\subsection{El problema que resuelve}

Sin Toku, una organización que cobra mensualmente (un gimnasio, una universidad, una aseguradora)
tiene que:

\begin{itemize}
  \item Emitir una boleta/factura manualmente
  \item Esperar que el cliente pague (transferencia, efectivo, tarjeta)
  \item Conciliar manualmente quién pagó y quién no
  \item Llamar a los morosos
  \item Todo esto tarda días/semanas y tiene alto costo operativo
\end{itemize}

Con Toku:
\begin{itemize}
  \item La organización carga los clientes y deudas vía API o CSV
  \item Toku cobra automáticamente en la fecha indicada
  \item El sistema notifica el resultado vía webhook
  \item La conciliación es automática
  \item Los reintentos son automáticos
\end{itemize}

\subsection{Por qué es técnicamente interesante}

\begin{tcolorbox}[colback=tokucyan!5,colframe=tokucyan]
\textbf{Toku no es un procesador de pagos} (no compite con Stripe o Kushki directamente).
Es un \textbf{orquestador} de cobros recurrentes que se conecta encima de los sistemas
bancarios de cada país. Eso es técnicamente complejo porque cada país tiene protocolos
distintos, regulaciones distintas y APIs bancarias distintas.
\end{tcolorbox}

\subsection{Su propuesta de valor en números}

\begin{center}
\begin{tabular}{lcc}
\toprule
\textbf{Métrica} & \textbf{Sin Toku} & \textbf{Con Toku} \\
\midrule
Costo por transacción & Alto (tarjeta ~2-3\%) & Bajo (débito ~0.5-1\%) \\
Tasa de pago exitoso & 70-80\% & Hasta 95\% \\
Tiempo de conciliación & Días & Automático \\
Reintentos manuales & Sí (costoso) & Automático \\
\bottomrule
\end{tabular}
\end{center}

% ─── SECCIÓN 2: DÉBITO DIRECTO ───────────────────────────────────────────────
\section{Débito Directo vs Cobro con Tarjeta}

\subsection{¿Qué es el débito directo?}

\begin{dominiobox}{Débito directo}
El cliente autoriza a la organización a cobrar \textbf{directamente de su cuenta bancaria}
mediante un \textbf{mandato} (también llamado \textit{autorización}, \textit{PAC} en Chile).
El cobro es pull: la organización "jala" el dinero, sin que el cliente haga nada.
\end{dominiobox}

\subsection{Ciclo de vida de un mandato de débito directo}

\begin{enumerate}
  \item \textbf{Autorización}: El cliente firma/acepta el mandato digitalmente o en papel
  \item \textbf{Registro}: El banco del cliente registra el mandato
  \item \textbf{Cobro}: La organización envía un lote de cobros al banco
  \item \textbf{Procesamiento}: El banco verifica fondos y débita
  \item \textbf{Respuesta}: El banco confirma éxito o reporta error (fondos insuficientes, cuenta cerrada)
  \item \textbf{Reintento}: Si falla, el sistema puede reintentar automáticamente
\end{enumerate}

\subsection{Comparación de métodos}

\begin{center}
\begin{tabular}{p{3cm}p{4cm}p{4cm}p{3cm}}
\toprule
\textbf{Aspecto} & \textbf{Débito directo} & \textbf{Tarjeta} & \textbf{Transferencia manual} \\
\midrule
Costo & \$\quad (bajo) & \$\$\$ (2-3\%) & \$\quad (bajo) \\
Fricción para cliente & Autorización única & Nada & Alta (hace la transferencia) \\
Tasa de éxito & Alta (95\%+) & Media (tarjeta vence) & Baja (olvida pagar) \\
Riesgo chargeback & Bajo & Alto & Ninguno \\
Velocidad liquidación & 1-3 días hábiles & Inmediato & Inmediato \\
Aplica para & Recurrente & Recurrente/único & Único \\
\bottomrule
\end{tabular}
\end{center}

% ─── SECCIÓN 3: SISTEMAS POR PAÍS ────────────────────────────────────────────
\section{Sistemas de Pago por País}

\begin{clbox}
\textbf{PAC} = Pago Automático en Cuenta (débito desde cuenta bancaria)\\
\textbf{PAT} = Pago Automático con Tarjeta (débito desde tarjeta de débito/crédito)

\vspace{6pt}
\textbf{Cómo funciona PAC en Chile:}
\begin{enumerate}
  \item Cliente firma mandato (físico o web)
  \item Organización registra el mandato en su banco (Banco A)
  \item El día de cobro, Banco A envía instrucción al Banco B del cliente
  \item SINACOFI (cámara de compensación) procesa las instrucciones en lotes
  \item Respuesta al día siguiente hábil (T+1)
\end{enumerate}

\textbf{Particularidades técnicas:}
\begin{itemize}
  \item El cobro es en \textbf{lotes}, no en tiempo real
  \item Hay cortes horarios (generalmente a las 14:00)
  \item Los errores llegan al día siguiente
  \item Número de cuenta + RUT del cliente son los identificadores
\end{itemize}
\end{clbox}

\vspace{8pt}

\begin{mxbox}
\textbf{CLABE} = Clave Bancaria Estandarizada (18 dígitos, identifica banco+sucursal+cuenta)\\
\textbf{SPEI} = Sistema de Pagos Electrónicos Interbancarios (transferencias en tiempo real)

\vspace{6pt}
\textbf{Débito en México:}
\begin{itemize}
  \item Se llama \textbf{CoDi} (Cobro Digital) o débito domiciliado
  \item El cliente autoriza vía su banco móvil
  \item Las instrucciones de cobro van vía CECOBAN (cámara de compensación)
  \item SPEI para transferencias inmediatas (24/7)
  \item La CLABE tiene checksum integrado (dígito verificador)
\end{itemize}

\textbf{Facturación electrónica (SAT):}\\
En México toda transacción comercial requiere CFDI (factura electrónica). Toku integra esto.
\end{mxbox}

\vspace{8pt}

\begin{brbox}
\textbf{PIX} = Sistema de pagos instantáneos del Banco Central de Brasil (BCB)

\vspace{6pt}
\textbf{Características únicas:}
\begin{itemize}
  \item Lanzado en 2020, 24/7/365, respuesta en segundos
  \item Identificadores: CPF/CNPJ, email, teléfono, o chave aleatória
  \item \textbf{Pix Automático} (2024): permite débito recurrente vía Pix
  \item Boleto: método tradicional, vence a los 3 días, pago en efectivo posible
  \item Débito automático: similar a PAC, requiere autorización con banco
\end{itemize}

\textbf{Por qué Brasil es complejo:}\\
Tiene 3 métodos distintos (Pix, Boleto, Débito Automático), cada uno con
regulación del BCB diferente. Toku tiene que manejar los 3.
\end{brbox}

% ─── SECCIÓN 4: FLUJOS COMPLETOS ─────────────────────────────────────────────
\section{Flujos Completos de Pago}

\subsection{Flujo 1: Cobro recurrente exitoso (el happy path)}

\begin{tcolorbox}[colback=tokucyan!5,colframe=tokucyan]
\begin{enumerate}
  \item \texttt{POST /customers} → cliente registrado en Toku
  \item \texttt{POST /subscriptions} → se crea agrupador con product\_id
  \item \texttt{POST /invoices} → se crea la deuda del mes
  \item \texttt{POST /transactions} → Toku intenta el cobro
  \item Toku envía instrucción al sistema bancario
  \item Sistema bancario responde: éxito
  \item Toku actualiza invoice → \texttt{estado = 'paid'}
  \item \texttt{webhook: invoice.paid} → tu sistema actualiza su BD
  \item \texttt{webhook: transaction.completed} → confirmación final
\end{enumerate}
\end{tcolorbox}

\subsection{Flujo 2: Cobro fallido con reintento automático}

\begin{tcolorbox}[colback=alertred!5,colframe=alertred]
\begin{enumerate}
  \item Mismo inicio que Flujo 1...
  \item Sistema bancario responde: error (fondos insuficientes)
  \item \texttt{webhook: invoice.failed} → tu sistema recibe la notificación
  \item Toku espera N horas (configurable)
  \item \texttt{POST /transactions} (intento 2) → nuevo intento de cobro
  \item Si falla 5 veces: invoice queda en \texttt{failed}, notificación final
  \item Tu sistema puede enviar email al cliente para que actualice su info
\end{enumerate}
\end{tcolorbox}

\subsection{Flujo 3: Cliente se inscribe en el portal}

\begin{tcolorbox}[colback=tipgreen!5,colframe=tipgreen]
\begin{enumerate}
  \item Tu backend genera checkout session: \texttt{POST /checkout\_sessions}
  \item Devuelve \texttt{session\_id}
  \item Frontend muestra el iframe de Toku con ese session\_id
  \item Cliente ingresa datos bancarios (PAC) o tarjeta
  \item Toku valida y registra el payment\_method
  \item \texttt{webhook: payment\_method.created} → tu sistema lo sabe
  \item Ahora puedes cobrarle automáticamente
\end{enumerate}
\end{tcolorbox}

% ─── SECCIÓN 5: IDEMPOTENCIA ─────────────────────────────────────────────────
\section{Idempotencia — El Concepto más Importante}

\begin{dominiobox}{Definición}
Una operación es \textbf{idempotente} si ejecutarla múltiples veces produce el mismo resultado
que ejecutarla una sola vez. En pagos esto es crítico: nunca puedes cobrar dos veces por error.
\end{dominiobox}

\subsection{¿Por qué pasa?}

\begin{itemize}
  \item Tu servidor envía la petición, el banco la procesa, pero la respuesta se pierde (timeout)
  \item Tu sistema no sabe si el cobro ocurrió → reintenta → doble cobro
  \item El webhook llega dos veces (redes no garantizan entrega única)
  \item Un worker de jobs ejecuta dos veces la misma tarea
\end{itemize}

\subsection{Soluciones prácticas}

\begin{lstlisting}[style=python]
# Patron 1: Idempotency Key en la API
import httpx
import uuid

async def crear_transaccion(invoice_id: str, monto: int) -> dict:
    idempotency_key = f"txn-{invoice_id}"  # determinista, no random!

    async with httpx.AsyncClient() as client:
        response = await client.post(
            "https://api.trytoku.com/transactions",
            headers={
                "Authorization": f"Bearer {API_KEY}",
                "Idempotency-Key": idempotency_key,  # Toku usa esto
            },
            json={"invoice_id": invoice_id, "amount": monto}
        )
    # Si se llama dos veces con la misma key, Toku devuelve el mismo resultado

# Patron 2: Chequear estado antes de actuar
async def procesar_invoice(invoice_id: str):
    # Primero verificar que no fue procesada ya
    invoice = await obtener_invoice(invoice_id)
    if invoice["estado"] != "pending":
        return  # Ya fue procesada, no hacer nada

    # Patron lock con PostgreSQL (SELECT FOR UPDATE)
    async with db.transaction():
        row = await db.fetchrow(
            "SELECT id FROM invoices WHERE id=$1 AND estado='pending' FOR UPDATE",
            invoice_id
        )
        if not row:
            return  # Otro proceso la tomó
        await cobrar_y_actualizar(invoice_id)

# Patron 3: Idempotencia en webhooks
async def handle_webhook(event: dict):
    event_id = event["id"]

    # Intentar insertar el event_id
    resultado = await db.execute(
        """INSERT INTO webhook_events (event_id, event_type, payload)
           VALUES ($1, $2, $3)
           ON CONFLICT (event_id) DO NOTHING""",
        event_id, event["type"], event
    )
    if resultado == "INSERT 0":
        return  # Duplicado, ignorar

    # Procesar el evento
    await dispatch_event(event)
\end{lstlisting}

% ─── SECCIÓN 6: CONCEPTOS CLAVE PARA LA ENTREVISTA ───────────────────────────
\section{Conceptos Clave que Debes Dominar}

\begin{center}
\begin{tabular}{p{3.5cm}p{10.5cm}}
\toprule
\textbf{Concepto} & \textbf{En el contexto de Toku} \\
\midrule
\textbf{Idempotencia} & Nunca cobrar dos veces. Usar idempotency keys en la API. Registrar event\_id en webhooks. \\
\textbf{Reintentos} & Lógica de backoff exponencial. Máximo N intentos. Notificar si todos fallan. \\
\textbf{Mandato/Autorización} & El consentimiento del cliente para cobrarle. Sin esto no se puede debitar. \\
\textbf{Conciliación} & Verificar que lo que Toku dice que cobró coincide con lo que tu sistema registra. \\
\textbf{Settlement} & Cuando Toku te deposita los fondos cobrados (puede ser D+1, D+2 según banco). \\
\textbf{Chargeback} & El cliente disputa el cobro con su banco. El banco te devuelve el cargo. \\
\textbf{Webhook} & Notificación que Toku te envía cuando algo cambia. Tu sistema reacciona a ellos. \\
\textbf{External ID} & Tu propio ID para sincronizar con tu sistema interno. Evita duplicados. \\
\textbf{Product ID} & Agrupa facturas del mismo tipo de servicio (ej: "arriendo enero"). \\
\textbf{Estado de invoice} & pending → paid / failed / void / partially\_paid. Máquina de estados. \\
\bottomrule
\end{tabular}
\end{center}

\subsection{Preguntas frecuentes en entrevistas de dominio}

\begin{tcolorbox}[colback=warnambel!5,colframe=warnambel,title=\textbf{Pregunta probable en entrevista}]
\textit{"¿Cómo manejarías el caso en que Toku te dice que la transacción fue exitosa,
pero tu base de datos falló al actualizar el estado de la factura?"}

\textbf{Respuesta esperada:}
\begin{itemize}
  \item Los webhooks son la red de seguridad: aunque tu UPDATE falle, el webhook llegará
  \item Diseñar el handler de webhooks para ser idempotente
  \item Usar transacciones en la BD: UPDATE invoice + INSERT en evento\_log como una sola transacción
  \item Tener un proceso de reconciliación periódica que compare tu BD con la API de Toku
\end{itemize}
\end{tcolorbox}

\begin{tcolorbox}[colback=warnambel!5,colframe=warnambel,title=\textbf{Pregunta probable en entrevista}]
\textit{"¿Cómo diseñarías el sistema de reintentos de cobro?"}

\textbf{Respuesta esperada:}
\begin{itemize}
  \item Backoff exponencial: esperar 1h, 4h, 24h entre intentos
  \item Máximo de intentos configurable (default 5)
  \item Diferenciar errores retriables (fondos insuficientes) de no retriables (cuenta cerrada)
  \item Notificar al cliente después de N intentos fallidos
  \item Job scheduler (Celery/BullMQ) con delay configurable
\end{itemize}
\end{tcolorbox}

\end{document}
